% page 393 du fac-simile (image p-392 du scan)
% bloc 165.1 x 233.3 mm ; marge gauche 36.9 mm
\pgc{15.240mm}{0.000mm}{
\l{\cel{52}385.}
\l{}
\l{\sou{moyeno}. To quo uzesas por atingar la skopo. F.}
\l{}
\l{\sou{mozaiko.}\cel{2}Laboruro ek peci inter-adjustita ek stono, petro,}
\l{\cel{3}emalio, glaso, ligno, diversa-kolora, di qui la asembluro}
\l{\cel{3}formacas desegnuri od ornamenti. - DEFIRS.}
\l{}
\l{\sou{muar}. (\sou{netr.}) I. (\sou{aludante ula speci di animali}). Subisar, ye}
\l{\cel{3}ula epoki determinita di lia vivo-duro, la falo o la rinovigo}
\l{\cel{3}di epidermo, pili, plumi, korni, e c. - II.(\sou{aludante la voco}}
\l{\cel{3}\sou{homala)}. Subisar ula altero lor puberesko. - EFIS.}
\l{}
\l{\sou{muaro}. Stofo di qua la reflekti esas kolor-chanjanta qui pro-}
\l{\cel{3}duktesas da la aplasto di la grani sub la cilindri inter-}
\l{\cel{3}rotacanta.- DEFIRS.}
\l{}
\l{\sou{mucilajo}. I. Substanco viskoza quan kontenas ula vejetanti. -}
\l{\cel{3}II. (\sou{farmaci}o) Liquido viskoza, ek la solvo di gumo en aquo.}
\l{\cel{3}- EFIS.}
\l{}
\l{\sou{muelar}. (\sou{trans}.)Triturar (grani) per pasigar inter grosa diski}
\l{\cel{3}qui rotacas, ordinare ek petro harda - la supra frotante}
\l{\cel{3}kontre la infra. - FIRS.}
\l{}
\l{\sou{muevo}. (\sou{zool.}) Mar-ucelo, palmipeda, long-ala. - L.\cel{1}\sou{larus}.-}
\l{\cel{3}DF.}
\l{}
\l{\sou{mufo}. I. Furajo, gaino vatizita, apertita ye la du extremaji,}
\l{\cel{3}en qua onu pozas sua manui por shirmar olci de la atmosfero}
\l{\cel{3}kolda. - II. (\sou{teknol}.)Cilindro ek gisfero aden qua onu ajust-}
\l{\cel{3}as la extremaji di du tubi, di du ax-arbori, quin onu volas}
\l{\cel{3}interjuntar. -DER.}
\l{}
\l{\sou{muflo}. I. (\sou{teknol}.)Tera vazo quan la kemiisti uzas por submisar}
\l{\cel{3}korpo a la efiko da fairo sen kontakto kun la flamo. -}
\l{\cel{3}II. Mikra forno quan onu pozas transverse en forno plu granda,}
\l{\cel{3}e qua recevas la materii destinita a kupelago. - DEFIS.}
\l{}
\l{\sou{muflono}. (\sou{zool.)}\cel{2}Sorto di mutono sovaja, ne-amansita (en Kor-}
\l{\cel{3}sika, Sardinia, Grekia). - L.\cel{1}\sou{ovis musimo.}- DEFIS.}
\l{}
\l{\sou{mujar.}\cel{2}(\sou{netrans}.)(\sou{aludante vento, ondi)}. Produktar bruiso ana-}
\l{\cel{3}loga a la klamo matida, nesonora e longa di la animali bova.}
\l{\cel{3}- FIS.}
\l{}
\l{\sou{muko}. (\sou{fiziol.)}\cel{2}Liquido viskoza quan sekrecas ula membrani. -}
\l{\cel{3}EFIS.}
\l{}
\l{\sou{mulo}. (\sou{zool.}) Produkturo maskula ek la kopulaco da asnulo kun}
\l{\cel{3}kavalino, o da kavalulo kun asnino. - EFIRS.}
\l{}
\l{mulato. Individuo qua naskis de la koito da pelblanko kun pel-}
\l{\cel{3}nigro. - DEFIRS}
}
